% GENERATED FILE. Edit canonical lessons, then run npm run generate:books.
% canonical-source-hash: fnv1a64:92f6743df9234b2a
% canonical-lessons: ZH-C129-anjing, ZH-C129-chao, ZH-W129-chao-noisy, ZH-C129-renao, ZH-C129-fangbian, ZH-C129-youyisi

\chapter{Quiet and Noisy}
\label{ch:zh-quiet-and-noisy}
\hlchaptermodality{129}

\begin{chapteropening}
\textbf{By the end of this chapter:} \emph{I can say a place is quiet, noisy, lively or handy, and that something is interesting.}

Complete the last lesson of chapter 129: I can say a place is quiet, noisy, lively or handy, and that something is interesting.
\end{chapteropening}

\section[ānjìng]{ānjìng — quiet}

\label{lesson:ZH-C129-anjing}



\begin{quote}
Before the new one: say the Mandarin for concert, then the Mandarin for guitar.
\end{quote}

\subsection*{You'll want to know: ānjìng}
\emph{ānjìng} — \textquotedblleft{}quiet\textquotedblright{}.

\begin{quote}
\emph{wǒ de fángjiān hěn ānjìng} — My room is quiet.
\end{quote}

Its characters come later; for now it is a word for the ear.

\begin{cousinweb}
\emph{ānjìng} is built from \emph{ān} 'peaceful' and \emph{jìng} 'calm', both inherited from Middle Chinese, \emph{'an} and \emph{dzjengX} (Guangyun, in Baxter's transcription); the dictionary glosses the second with the first.
\end{cousinweb}

\begin{grammarlens}[title={What you've built}]
One more word for describing things.
\end{grammarlens}

\subsection*{Your turn}
\begin{itemize}
\raggedright
  \item \emph{Say it:} \emph{ānjìng}
  \item \emph{Say it:} \emph{ānjìng}, once more
  \item \emph{Say it:} \emph{wǒ de fángjiān hěn ānjìng} — My room is quiet.
  \item \emph{Recall:} say the Mandarin for concert, then the Mandarin for guitar, then say \emph{ānjìng} again
\end{itemize}

\begin{tcolorbox}[breakable,colback=teal!4,colframe=teal!35!black,title={Before you move on}]
What does ānjìng mean? (\textquotedblleft{}Quiet\textquotedblright{}.) Say it once more.
\end{tcolorbox}
\section[\texorpdfstring{\zh{吵} (chǎo)}{吵 (chǎo)}]{\zh{吵} (chǎo) — noisy}

\label{lesson:ZH-C129-chao}



\begin{quote}
Before the new one: say the Mandarin for guitar, then the Mandarin for quiet.
\end{quote}

\subsection*{You'll want to know: \zh{吵}}
\textbf{\zh{吵}} — \emph{chǎo} — \textquotedblleft{}noisy\textquotedblright{}.

\begin{quote}
\textbf{\zh{他的音乐太吵了}} \emph{tā de yīnyuè tài chǎo le} — His music is far too loud.
\end{quote}

The new character, in pieces: \zh{吵} — \zh{口} on the left and \zh{少} on the right. The next lesson puts it in your hand.

\begin{cousinweb}
\textbf{\zh{吵}} (\emph{chǎo}) is inherited from Middle Chinese \emph{tsrhaewX} (Guangyun, in Baxter's transcription), glossed 'sound', a different word from the \emph{chǎo} of stir-frying, which has the same old form. The dictionary notes that the character first had another reading, \emph{mjiewX}, glossed 'the call of a pheasant'.
\end{cousinweb}

\begin{grammarlens}[title={What you've built}]
One more word for describing things.
\end{grammarlens}

\subsection*{Your turn}
\begin{itemize}
\raggedright
  \item \emph{Say it:} \emph{chǎo}
  \item \emph{Say it:} \emph{chǎo}, once more
  \item \emph{Say it:} \emph{tā de yīnyuè tài chǎo le} — His music is far too loud.
  \item \emph{Recall:} say the Mandarin for guitar, then the Mandarin for quiet, then say \emph{chǎo} again
\end{itemize}

\begin{tcolorbox}[breakable,colback=teal!4,colframe=teal!35!black,title={Before you move on}]
What does \zh{吵} mean? (\textquotedblleft{}Noisy\textquotedblright{}.) Say it once more.
\end{tcolorbox}
\section[\texorpdfstring{\zh{吵} (chǎo)}{吵 (chǎo)}]{\zh{吵} — \zh{口}, then \zh{少}}

\label{lesson:ZH-W129-chao-noisy}



\begin{quote}
The word first, then the character inside it.

\begin{itemize}
\raggedright
  \item \emph{Recall:} say \emph{chǎo} — \textbf{R1}, one lesson back, before it decays
\end{itemize}

\begin{quote}
\zh{口} \zh{少}
\end{quote}

\emph{Read it:} each one aloud, with nothing said first — \emph{kǒu}, \emph{shǎo}
\end{quote}

\subsection*{Script — \zh{口}, then \zh{少}}
\hlblockfigure{\detokenize{figures/ZH-W129-chao-noisy-filmstrip.pdf}}{How \zh{吵} is written, stroke by stroke}

\begin{quote}
\zh{吵}
\end{quote}

\textbf{Two pieces, left before right}, and both are old friends.

\textbf{\zh{口}, three strokes}, small and narrow on the left: the left side, the top that turns down the right side, and the bar that closes the bottom.

\textbf{\zh{少}, four strokes}, on the right: the vertical down the middle, the short stroke on its left, the dot on its right, then the long left-falling stroke that sweeps under them.

3 + 4 = seven.

\subsection*{Your turn}
\begin{itemize}
\raggedright
  \item \emph{Write it:} \zh{吵}, all seven, with the model in front of you
  \item \emph{Cover:} write it again from memory, and count the pieces, not the strokes
\end{itemize}

\begin{tcolorbox}[breakable,colback=teal!4,colframe=teal!35!black,title={Before you move on}]
What are the two pieces of \zh{吵}? (\textbf{\zh{口} and \zh{少}}.) Which side is narrow? (\textbf{The left, \zh{口}}.) What does the last stroke do? (\textbf{It sweeps down to the left under \zh{少}}.)

Source: \href{https://raw.githubusercontent.com/chanind/hanzi-writer-data/68d10a4b21150cae5e1ebbd223eed289cf32d90c/data/\%E5\%90\%B5.json}{Hanzi Writer Data \zh{吵}, pinned stroke paths and medians}.
\end{tcolorbox}
\section[rènao]{rènao — lively; bustling}

\label{lesson:ZH-C129-renao}



\begin{quote}
Before the new one: say the Mandarin for quiet, then the Mandarin for noisy.
\end{quote}

\subsection*{You'll want to know: rènao}
\emph{rènao} — \textquotedblleft{}lively; bustling\textquotedblright{}.

\begin{quote}
\emph{Běijīng hěn rènao} — Beijing is a lively place.
\end{quote}

Its characters come later; for now it is a word for the ear.

\begin{cousinweb}
\emph{rènao} is built from \emph{rè} 'hot' and \emph{nao}, both inherited from Middle Chinese, \emph{nyet} and \emph{nraewH} (Guangyun, in Baxter's transcription); the dictionary glosses the second 'not quiet' and 'to stir up'.
\end{cousinweb}

\begin{grammarlens}[title={What you've built}]
One more word for describing things.
\end{grammarlens}

\subsection*{Your turn}
\begin{itemize}
\raggedright
  \item \emph{Say it:} \emph{rènao}
  \item \emph{Say it:} \emph{rènao}, once more
  \item \emph{Say it:} \emph{Běijīng hěn rènao} — Beijing is a lively place.
  \item \emph{Recall:} say the Mandarin for quiet, then the Mandarin for noisy, then say \emph{rènao} again
\end{itemize}

\begin{tcolorbox}[breakable,colback=teal!4,colframe=teal!35!black,title={Before you move on}]
What does rènao mean? (\textquotedblleft{}Lively; bustling\textquotedblright{}.) Say it once more.
\end{tcolorbox}
\section[fāngbiàn]{fāngbiàn — convenient; handy}

\label{lesson:ZH-C129-fangbian}



\begin{quote}
Before the new one: say the Mandarin for noisy, then the Mandarin for bustling.
\end{quote}

\subsection*{You'll want to know: fāngbiàn}
\emph{fāngbiàn} — \textquotedblleft{}convenient; handy\textquotedblright{}.

\begin{quote}
\emph{zhèr mǎi dōngxi hěn fāngbiàn} — Shopping here is easy.
\end{quote}

Its characters come later; for now it is a word for the ear.

\begin{cousinweb}
\emph{fāngbiàn} is built from \emph{fāng} and \emph{biàn}, both inherited from Middle Chinese, \emph{pjang} and \emph{bjienH} (Guangyun, in Baxter's transcription). The dictionary glosses the first 'the four quarters' and 'a method', and the second 'advantageous'.
\end{cousinweb}

\begin{grammarlens}[title={What you've built}]
One more word for describing things.
\end{grammarlens}

\subsection*{Your turn}
\begin{itemize}
\raggedright
  \item \emph{Say it:} \emph{fāngbiàn}
  \item \emph{Say it:} \emph{fāngbiàn}, once more
  \item \emph{Say it:} \emph{zhèr mǎi dōngxi hěn fāngbiàn} — Shopping here is easy.
  \item \emph{Recall:} say the Mandarin for noisy, then the Mandarin for bustling, then say \emph{fāngbiàn} again
\end{itemize}

\begin{tcolorbox}[breakable,colback=teal!4,colframe=teal!35!black,title={Before you move on}]
What does fāngbiàn mean? (\textquotedblleft{}Convenient; handy\textquotedblright{}.) Say it once more.
\end{tcolorbox}
\section[yǒuyìsi]{yǒuyìsi — interesting}

\label{lesson:ZH-C129-youyisi}



\begin{quote}
Before the new one: say the Mandarin for bustling, then the Mandarin for convenient.
\end{quote}

\subsection*{You'll want to know: yǒuyìsi}
\emph{yǒuyìsi} — \textquotedblleft{}interesting\textquotedblright{}.

\begin{quote}
\emph{xué hànyǔ hěn yǒuyìsi} — Learning Chinese is interesting.
\end{quote}

Its characters come later; for now it is a word for the ear.

\begin{cousinweb}
\emph{yǒuyìsi} is built from \zh{有} 'to have' and \emph{yìsi} 'meaning', the \emph{yìsi} of \emph{bù hǎoyìsi}: to have meaning. The parts are inherited from Middle Chinese, \emph{hjuwX}, \emph{'iH} and \emph{si} (Guangyun, in Baxter's transcription).
\end{cousinweb}

\begin{grammarlens}[title={What you've built}]
One more word for describing things.
\end{grammarlens}

\subsection*{Your turn}
\begin{itemize}
\raggedright
  \item \emph{Say it:} \emph{yǒuyìsi}
  \item \emph{Say it:} \emph{yǒuyìsi}, once more
  \item \emph{Say it:} \emph{xué hànyǔ hěn yǒuyìsi} — Learning Chinese is interesting.
  \item \emph{Recall:} say the Mandarin for bustling, then the Mandarin for convenient, then say \emph{yǒuyìsi} again
\end{itemize}

\begin{tcolorbox}[breakable,colback=teal!4,colframe=teal!35!black,title={Before you move on}]
What does yǒuyìsi mean? (\textquotedblleft{}Interesting\textquotedblright{}.) Say it once more.
\end{tcolorbox}
